\section{The prospectively fixed R124 inference family}
\label{supp:inference124}
The following inference is conditional on the published scientific source,
the actual fixed actors, sound numerical enclosures, and independent uniform
entropy bits. The population law is the declared two-coordinate 40-bit
dyadic law for initial states. The innovations inside each Bellman operator
retain their continuous law. All 16 initial states are independent units;
different methods, preferences, candidate attempts and timing repetitions
on one state supply no additional independent observations.

\begin{proposition}[Simultaneous original-policy inference]
\label{prop:inference124}
Put $n=16$, $\tau=1/256$, and $t=2^{-24}$. Consider the fixed 116
policy cells and 290 signed comparisons specified in
\texttt{REPORTING\_SPEC124.md}. In every cell let $G(X)$ be the actual
gain relative to its own qualified incumbent. Assume that the service
either retains the incumbent or verifies a strictly positive gain, and
that both policies have cost between the original optimum and that
optimum plus $\tau$. Then $0\le G\le\tau$.
If the numerical sample mean is enclosed by $[\bar\ell,\bar u]$,
the interval
\[
 [\max(0,\bar\ell-5\tau/8),
       \min(\tau,\bar u+5\tau/8)]
\]
contains the population mean except on a joint event of probability at
most $1/100$ over all 116 cells.

In each cell also count states on which the verified gain is at least
$t$. Let $k$ be this count and let $p_k^-$ be a lower one-sided
Clopper--Pearson probability bound with failure $1/11600$.
The lower mean-gain bound $t p_k^-$ holds simultaneously for the 116
cells except on an event of probability at most $1/100$.

For each fixed signed comparison, positive means the first policy has
lower original cost. The statewise difference belongs to
$[-\tau,\tau]$. Its numerical sample-mean interval enlarged in both
directions by $5\tau/4$, then intersected with that support, contains
its population mean except on a joint event of probability at most
$3/100$ over all 290 comparisons. A method compared with itself has
exact population difference zero. Observed equality on the sampled
states alone establishes only a sample identity.

The intersection of the two valid lower bounds for each nonnegative
gain, and all signed intervals, therefore has simultaneous coverage
at least $19/20$.
\end{proposition}

\begin{proof}
An accepted change has nonnegative actual gain by its exact certificate;
retention has zero gain. Since the incumbent is no more than $\tau$
above the common original optimum and the candidate cannot fall below
that optimum, the gain is at most $\tau$. Two qualified policies are
both in that same interval of length $\tau$, proving the signed support.

Hoeffding's two-sided inequality for independent observations in an
interval of length $R$ has failure at most
$2\exp(-2nr^2/R^2)$ at radius $r$. For $R=\tau$ and
$r=5\tau/8$, and for $R=2\tau$ and $r=5\tau/4$, the exponent is
$25/2$. Exact rational arithmetic verifies
\[
 \sum_{j=0}^{30}\frac{(25/2)^j}{j!}>23200.
\]
Consequently $232\exp(-25/2)<1/100$ and
$580\exp(-25/2)<3/100$. These are union bounds over the fixed
families; dependence across methods on the same state does not affect
them. Replacing the exact sample mean by a valid enclosing interval
can only enlarge the resulting confidence interval.

The threshold event is a deterministic measurable function of the
fixed service and the independent state. Its count is binomial with
parameters $n$ and its population event probability $p$. At $k=0$
set $p_k^-=0$. Otherwise the lower endpoint is the solution of
\[
 \sum_{j=k}^{n}{n\choose j}p^j(1-p)^{n-j}=1/11600.
\]
Binomial tail inversion gives the stated one-sided coverage. The
implementation takes the lower endpoint after 48 exact rational
bisections, so numerical rounding cannot raise this lower bound.
On the threshold event $G\ge t$, hence
$\mathbb E G\ge t p\ge t p_k^-$. A union bound over 116 such
events costs at most $1/100$. The sum of the three family failure
bounds is $1/20$. All claims are conditional on the stated numerical
and sampling premises; none uses timing repeats as population data.
\end{proof}

The deliberately small sample has a practical limitation already visible
before execution: the signed radius exceeds $\tau$, so these broad
distribution-free signed intervals cannot certify a nonzero ordering.
Their purpose is a complete prospective account and detection of any
incorrect transfer of older population results, not a promised ordering.
Positive own-incumbent mean-gain lower bounds can still follow from the
threshold-event route. Neither route identifies field computation prices
or a population adoption decision.

For completeness, if $q$ is $\mu$-strongly convex on the feasible
interval and $a^*$ is its constrained minimizer, its variational
inequality and strong convexity give
\[
 q(a)-q(a^*)\ge\frac{\mu}{2}|a-a^*|^2.
\]
Since $q(a)\le J(a)$, any valid upper bound $r$ on
$J(a)-q(a^*)$ implies $|a-a^*|\le\sqrt{2r/\mu}$.
The reported root-action RMS bound uses $\mu=2$, averages the valid
regret bounds first, and rounds the square root upward. It is an upper
bound on distance to the optimal-future root action, not an observed
exact action error or a statement about the derivative of the deployed
future-policy error.
